\documentclass[12pt]{article}
\usepackage[utf8]{inputenc}
\usepackage[spanish]{babel}
\usepackage{amsmath, amsfonts, amssymb, mathrsfs}
\usepackage{geometry}
\usepackage{fancyhdr}
\usepackage{parskip}
\usepackage{float}
\usepackage{hyperref}
\usepackage{tikz}
\usetikzlibrary{arrows.meta,positioning}
\AtBeginEnvironment{tikzpicture}{\shorthandoff{<>}}

\geometry{a4paper, margin=2cm}

\pagestyle{fancy}
\fancyhf{}
\lhead{Astroestadística }
\chead{Problemset 4}


\begin{document}
\section{Problema 1}
\subsection{Enunciado}

\textbf{Consistencia de un estimador.} Un estimador $\hat{\theta}$ se dice que es \textbf{consistente} si para cualquier $\epsilon > 0$,

$$P(|\hat{\theta} - \theta| \geq \epsilon) \to 0 \quad \text{cuando } n \to \infty.$$

Es decir, $\hat{\theta}$ es consistente si, a medida que el tamaño de la muestra crece, es cada vez menos probable que $\hat{\theta}$ esté a una distancia mayor que $\epsilon$ del valor verdadero de $\theta$.

Demuestre que $\overline{X}$ es un estimador consistente de $\mu$ cuando $\sigma^2 < \infty$ utilizando la desigualdad de Chebyshev. [Sugerencia: La desigualdad puede escribirse de la forma

$$P(|Y - \mu_Y| \geq \epsilon) \leq \frac{\sigma_Y^2}{\epsilon^2}.$$

Ahora identifique $Y$ con $\overline{X}$.]

\subsection{Solución}

Para demostrar que $\overline{X}$ es un estimador consistente de $\mu$, debemos probar que para cualquier $\epsilon > 0$:

$$P(|\overline{X} - \mu| \geq \epsilon) \to 0 \quad \text{cuando } n \to \infty.$$

\textbf{Paso 1: Identificar los parámetros de $\overline{X}$}

Sabemos que para la media muestral $\overline{X}$:
\begin{itemize}
    \item Valor esperado: $E(\overline{X}) = \mu_{\overline{X}} = \mu$
    \item Varianza: $Var(\overline{X}) = \sigma^2_{\overline{X}} = \frac{\sigma^2}{n}$
\end{itemize}

\textbf{Paso 2: Aplicar la desigualdad de Chebyshev}

La desigualdad de Chebyshev establece que para cualquier variable aleatoria $Y$ con media $\mu_Y$ y varianza $\sigma^2_Y$:

$$P(|Y - \mu_Y| \geq \epsilon) \leq \frac{\sigma^2_Y}{\epsilon^2}$$

Identificando $Y = \overline{X}$, tenemos:

$$P(|\overline{X} - \mu| \geq \epsilon) \leq \frac{\sigma^2_{\overline{X}}}{\epsilon^2} = \frac{\sigma^2/n}{\epsilon^2} = \frac{\sigma^2}{n\epsilon^2}$$

\textbf{Paso 3: Tomar el límite cuando $n \to \infty$}

Ahora consideramos el límite cuando el tamaño de muestra tiende a infinito:

$$\lim_{n \to \infty} P(|\overline{X} - \mu| \geq \epsilon) \leq \lim_{n \to \infty} \frac{\sigma^2}{n\epsilon^2}$$

Dado que $\sigma^2 < \infty$ (varianza finita) y $\epsilon > 0$ es constante:

$$\lim_{n \to \infty} \frac{\sigma^2}{n\epsilon^2} = 0$$

Por lo tanto:

$$\lim_{n \to \infty} P(|\overline{X} - \mu| \geq \epsilon) = 0$$

\textbf{Conclusión:}

Hemos demostrado que $\overline{X}$ satisface la definición de consistencia. Es decir, $\overline{X}$ es un estimador consistente de $\mu$ cuando $\sigma^2 < \infty$.

Esto tiene sentido intuitivo: a medida que aumentamos el tamaño de la muestra, la varianza de $\overline{X}$ disminuye ($\sigma^2/n$), haciendo que $\overline{X}$ se concentre cada vez más alrededor del verdadero valor $\mu$.

\section{Problema 2}

\textbf{Edades de cúmulos estelares.} La siguiente información muestra la edad (Gyr) de algunos cúmulos estelares masivos observados en La galaxia, calculados con modelos de evolución estelar.

\begin{center}
    \begin{tabular}{ccccccc}
    9.7 & 7.7 & 7.4 & 9.7 & 7.8 & 11.8 & 8.2 \\
    11.3 & 8.7 & 6.8 & 5.9 & 7.2 & 7.3 & 6.3 \\
    9.0 & 7.9 & 6.3 & 7.6 & 10.7 & 6.5 & 7.0 \\
    8.1 & 6.8 & 7.0 & 7.8 & 7.7 & 11.6 & \\
    \end{tabular}
\end{center}

\subsection{Solución}

Primero, ordenamos los datos de menor a mayor para facilitar los cálculos. El tamaño de la muestra es $n=27$.
\begin{center}
    5.9, 6.3, 6.3, 6.5, 6.8, 6.8, 7.0, 7.0, 7.2, 7.3, 7.4, 7.6, 7.7, 7.7, 7.8, 7.8, 7.9, 8.1, 8.2, 8.7, 9.0, 9.7, 9.7, 10.7, 11.3, 11.6, 11.8
\end{center}

\subsubsection*{a) Establezca un estimador para valor medio de la edad de todos los cúmulos estelares masivos en la galaxia y estime.}

El estimador natural para la media poblacional $\mu$ es la media muestral $\bar{X}$.
La estimación puntual se calcula como:
$$ \bar{x} = \frac{1}{n} \sum_{i=1}^{n} x_i $$
Sumando todos los valores observados:
$$ \sum x_i = 219.8 $$
Por lo tanto:
$$ \bar{x} = \frac{219.8}{27} \approx 8.14 \text{ Gyr} $$
La edad promedio estimada de los cúmulos estelares masivos es de aproximadamente 8.14 Gyr.

\subsubsection*{b) Hacer una estimación puntual del valor de edad que separa al 50\% de los cúmulos más jóvenes de la galaxia del 50\% de los cúmulos más viejos. Establezca el estimador que va a usar.}

El valor que separa el 50\% inferior del 50\% superior es la mediana. El estimador es la mediana muestral $\tilde{X}$.
Dado que $n=27$ es impar, la mediana es el valor que ocupa la posición central, es decir, la posición $(n+1)/2 = (27+1)/2 = 14$.
Observando los datos ordenados, el valor en la posición 14 es 7.7.
$$ \tilde{x} = 7.7 \text{ Gyr} $$

\subsubsection*{c) Calcular e interpretar un estimado puntual de la desviación estándar $\sigma$ de toda la población. ¿Qué estimador va a usar?}

El estimador insesgado para la desviación estándar poblacional $\sigma$ es la desviación estándar muestral $S$, definida como:
$$ S = \sqrt{\frac{1}{n-1} \sum_{i=1}^{n} (x_i - \bar{x})^2} $$
Primero calculamos la suma de cuadrados de las desviaciones o usamos la fórmula computacional $S_{xx} = \sum x_i^2 - \frac{(\sum x_i)^2}{n}$.
$$ \sum x_i^2 = 9.7^2 + 7.7^2 + \dots + 11.6^2 = 1860.94 $$
$$ S_{xx} = 1860.94 - \frac{(219.8)^2}{27} = 1860.94 - 1789.335 \approx 71.605 $$
Ahora calculamos la varianza muestral $s^2$:
$$ s^2 = \frac{71.605}{27-1} = \frac{71.605}{26} \approx 2.754 $$
Finalmente, la desviación estándar muestral $s$:
$$ s = \sqrt{2.754} \approx 1.66 \text{ Gyr} $$
\textbf{Interpretación:} En promedio, la edad de un cúmulo estelar masivo se desvía aproximadamente 1.66 Gyr de la edad media de la población.

\subsubsection*{d) Hacer una estimación puntual de la proporción de todos los cúmulos cuya edad excede los 10 Gyr}

El estimador para una proporción poblacional $p$ es la proporción muestral $\hat{P}$.
Contamos el número de observaciones mayores a 10: $\{11.8, 11.3, 10.7, 11.6\}$. Hay 4 observaciones.
$$ \hat{p} = \frac{4}{27} \approx 0.148 $$
Se estima que aproximadamente el 14.8\% de los cúmulos tienen una edad mayor a 10 Gyr.

\subsubsection*{e) Hacer una estimación puntual del coeficiente de variación $\sigma/\mu$ ¿Qué estimador va a usar?}

El estimador natural para el coeficiente de variación poblacional es el coeficiente de variación muestral, dado por $S/\bar{X}$.
Usando los valores calculados en los incisos anteriores:
$$ \frac{s}{\bar{x}} = \frac{1.66}{8.14} \approx 0.204 $$
El coeficiente de variación estimado es 0.204 (o 20.4\%), lo que indica que la desviación estándar es aproximadamente el 20\% del valor de la media.


\section{Problema 3}

\textbf{Enunciado:}

¿Con o sin azúcar? De un número $n_1$ de hombres seleccionados aleatoriamente, $X_1$ toman tinto con azúcar, mientras de $n_2$ mujeres, $X_2$ toman tinto con azúcar. Sean $p_1$ y $p_2$ las probabilidades de que un hombre y una mujer, respectivamente, seleccionados aleatoriamente tomen tinto con azúcar.

\begin{itemize}
    \item[a)] Muestre que $(X_1/n_1) - (X_2/n_2)$ es un estimador no sesgado para $p_1 - p_2$.
    \item[b)] ¿Cuál es el error estándar de este estimador?
    \item[c)] ¿Cómo usaría los valores observados $x_1$ y $x_2$ para estimar el error estándar de su estimador?
    \item[d)] Si $n_1 = n_2 = 200$, $x_1 = 127$ y $x_2 = 176$, use el estimador del literal (a) para obtener un estimado de $p_1 - p_2$.
    \item[e)] Use los datos anteriores para estimar el error estándar del estimador.
\end{itemize}

\subsection{Solución}

Este problema trata sobre la inferencia estadística para la diferencia entre dos proporciones de población. Los conceptos de estimadores puntuales, insesgamiento y error estándar se discuten en el Capítulo 6, mientras que la inferencia específica para dos muestras se aborda en el Capítulo 9, Sección 9.4 del libro de texto.

\subsubsection*{a) Demostración de que el estimador es no sesgado}

Para demostrar que un estimador es no sesgado, debemos probar que su valor esperado es igual al parámetro que se desea estimar. En este caso, el parámetro es la diferencia de proporciones poblacionales, $p_1 - p_2$.

El estimador propuesto es la diferencia de las proporciones muestrales:
$$ \hat{p}_1 - \hat{p}_2 = \frac{X_1}{n_1} - \frac{X_2}{n_2} $$

Donde $X_1$ y $X_2$ son variables aleatorias binomiales con parámetros $(n_1, p_1)$ y $(n_2, p_2)$ respectivamente. Sabemos que el valor esperado de una variable binomial es $E[X] = np$. Por lo tanto:
$$ E[X_1] = n_1 p_1 \quad \text{y} \quad E[X_2] = n_2 p_2 $$

Aplicando la propiedad de linealidad del valor esperado ($E[aX + bY] = aE[X] + bE[Y]$), calculamos el valor esperado del estimador:

$$ E\left[ \frac{X_1}{n_1} - \frac{X_2}{n_2} \right] = E\left[ \frac{X_1}{n_1} \right] - E\left[ \frac{X_2}{n_2} \right] $$
$$ = \frac{1}{n_1} E[X_1] - \frac{1}{n_2} E[X_2] $$
$$ = \frac{1}{n_1} (n_1 p_1) - \frac{1}{n_2} (n_2 p_2) $$
$$ = p_1 - p_2 $$

Dado que $E[\hat{p}_1 - \hat{p}_2] = p_1 - p_2$, hemos demostrado que $(X_1/n_1) - (X_2/n_2)$ es un estimador no sesgado para la diferencia de proporciones $p_1 - p_2$.

\subsubsection*{b) Error estándar del estimador}

El error estándar de un estimador es su desviación estándar. Para encontrarla, primero calculamos la varianza del estimador. Asumiendo que las muestras de hombres y mujeres son independientes, la varianza de la diferencia es la suma de las varianzas:

$$ V\left( \frac{X_1}{n_1} - \frac{X_2}{n_2} \right) = V\left( \frac{X_1}{n_1} \right) + V\left( \frac{X_2}{n_2} \right) $$

Usando la propiedad de la varianza $V(aX) = a^2 V(X)$ y sabiendo que la varianza de una variable binomial es $V(X) = np(1-p)$:

$$ V\left( \frac{X_1}{n_1} \right) = \frac{1}{n_1^2} V(X_1) = \frac{1}{n_1^2} n_1 p_1 (1 - p_1) = \frac{p_1 (1 - p_1)}{n_1} $$
$$ V\left( \frac{X_2}{n_2} \right) = \frac{1}{n_2^2} V(X_2) = \frac{1}{n_2^2} n_2 p_2 (1 - p_2) = \frac{p_2 (1 - p_2)}{n_2} $$

Por lo tanto, la varianza del estimador es:
$$ V(\hat{p}_1 - \hat{p}_2) = \frac{p_1 (1 - p_1)}{n_1} + \frac{p_2 (1 - p_2)}{n_2} $$

El error estándar ($\sigma_{\hat{p}_1 - \hat{p}_2}$) es la raíz cuadrada de la varianza:
$$ \sigma_{\hat{p}_1 - \hat{p}_2} = \sqrt{\frac{p_1 (1 - p_1)}{n_1} + \frac{p_2 (1 - p_2)}{n_2}} $$

\subsubsection*{c) Estimación del error estándar con valores observados}

En la práctica, los valores de $p_1$ y $p_2$ son desconocidos. Para estimar el error estándar, sustituimos las proporciones poblacionales $p_1$ y $p_2$ en la fórmula del inciso (b) con sus estimadores puntuales muestrales $\hat{p}_1 = x_1/n_1$ y $\hat{p}_2 = x_2/n_2$.

El error estándar estimado ($\hat{\sigma}_{\hat{p}_1 - \hat{p}_2}$) es:
$$ \hat{\sigma}_{\hat{p}_1 - \hat{p}_2} = \sqrt{\frac{\hat{p}_1 (1 - \hat{p}_1)}{n_1} + \frac{\hat{p}_2 (1 - \hat{p}_2)}{n_2}} $$
donde $\hat{p}_1 = \frac{x_1}{n_1}$ y $\hat{p}_2 = \frac{x_2}{n_2}$.

\subsubsection*{d) Estimación puntual de $p_1 - p_2$}

Con los datos proporcionados:
\begin{itemize}
    \item $n_1 = 200, x_1 = 127$
    \item $n_2 = 200, x_2 = 176$
\end{itemize}

Primero calculamos las proporciones muestrales:
$$ \hat{p}_1 = \frac{127}{200} = 0.635 $$
$$ \hat{p}_2 = \frac{176}{200} = 0.880 $$

El estimado de la diferencia $p_1 - p_2$ es simplemente la diferencia de las proporciones muestrales:
$$ \hat{p}_1 - \hat{p}_2 = 0.635 - 0.880 = -0.245 $$

Esto sugiere que la proporción de mujeres que toman tinto con azúcar es 0.245 (o 24.5\%) mayor que la de los hombres en la población.

\subsubsection*{e) Estimación del error estándar}

Usamos la fórmula desarrollada en el inciso (c) y los valores calculados en el inciso (d):

$$ \hat{\sigma}_{\hat{p}_1 - \hat{p}_2} = \sqrt{\frac{0.635 (1 - 0.635)}{200} + \frac{0.880 (1 - 0.880)}{200}} $$

Calculamos cada término dentro de la raíz:
1.  Para los hombres: $\frac{0.635 \times 0.365}{200} = \frac{0.231775}{200} \approx 0.0011589$
2.  Para las mujeres: $\frac{0.880 \times 0.120}{200} = \frac{0.1056}{200} = 0.0005280$

Sumamos los términos:
$$ 0.0011589 + 0.0005280 = 0.0016869 $$

Finalmente, tomamos la raíz cuadrada:
$$ \hat{\sigma}_{\hat{p}_1 - \hat{p}_2} = \sqrt{0.0016869} \approx 0.04107 $$

El error estándar estimado de la diferencia de proporciones es aproximadamente 0.0411.


\section{Problema 4}

\subsection{Enunciado}

\textbf{4. Distribución de Weibull.}

a) Consulte qué es la distribución de Weibull, escriba su pdf (función de densidad de probabilidad) y grafique.

b) Sea $X$ una variable aleatoria con distribución de Weibull con parámetros $\alpha$ y $\beta$ de modo que:
$$ E(X) = \beta\Gamma(1 + 1/\alpha) $$
$$ V(X) = \beta^2\{\Gamma(1 + 2/\alpha) - (\Gamma(1 + 1/\alpha))^2\} $$

Basándose en una muestra $X_1, ..., X_n$, utilice el método de los momentos para escribir las ecuaciones de los estimadores de $\alpha$ y $\beta$. Muestre que, una vez obtenido el estimado de $\alpha$, el estimado de $\beta$ puede obtenerse a partir de la función gamma y que el estimado de $\alpha$ se puede obtener de una ecuación que implica la función gamma.

Ahora, sean $n = 20$, $\bar{x} = 28.0$ y $\sum x_i^2 = 16500$, estime los parámetros $\beta$ y $\alpha$. [Ayuda: $[\Gamma(1.2)]^2 /\Gamma(1.4) = 0.95$]

\subsection{Solución}

\textbf{a) Distribución de Weibull y su gráfica}

De acuerdo con la \textbf{Sección 4.5 (Otras distribuciones continuas)} del libro de texto de Devore, la distribución de Weibull es una familia de distribuciones de probabilidad continua utilizada frecuentemente en análisis de confiabilidad y supervivencia.

La función de densidad de probabilidad (pdf) para una variable aleatoria $X$ con distribución de Weibull con parámetros de forma $\alpha > 0$ y de escala $\beta > 0$ se define como:

$$ f(x; \alpha, \beta) = \begin{cases} 
\frac{\alpha}{\beta^\alpha} x^{\alpha-1} e^{-(x/\beta)^\alpha} & x \geq 0 \\
0 & \text{en otro caso}
\end{cases} $$

\textit{Gráfica:} La forma de la gráfica depende del parámetro de forma $\alpha$.
\begin{itemize}
    \item Si $0 < \alpha < 1$, la función es decreciente y tiende a infinito cuando $x \to 0$.
    \item Si $\alpha = 1$, la distribución se reduce a una distribución exponencial.
    \item Si $\alpha > 1$, la función es unimodal (tiene forma de campana asimétrica), comenzando en 0, aumentando hasta un máximo y luego decayendo hacia 0.
\end{itemize}

\textbf{b) Estimación por el Método de los Momentos}

De acuerdo con la \textbf{Sección 6.2 (Métodos de estimación puntual)}, el método de los momentos consiste en igualar los momentos de la población con los momentos muestrales correspondientes.

Los momentos de la población dados son:
1. Primer momento (Media): $E(X) = \beta\Gamma(1 + 1/\alpha)$
2. Segundo momento central (Varianza): $V(X) = \beta^2\{\Gamma(1 + 2/\alpha) - (\Gamma(1 + 1/\alpha))^2\}$

Para aplicar el método, necesitamos el segundo momento raw de la población, $E(X^2)$. Sabemos que $V(X) = E(X^2) - [E(X)]^2$, por lo tanto:
$$ E(X^2) = V(X) + [E(X)]^2 $$
Sustituyendo las expresiones dadas:
$$ E(X^2) = \beta^2\{\Gamma(1 + 2/\alpha) - (\Gamma(1 + 1/\alpha))^2\} + [\beta\Gamma(1 + 1/\alpha)]^2 $$
$$ E(X^2) = \beta^2\Gamma(1 + 2/\alpha) - \beta^2(\Gamma(1 + 1/\alpha))^2 + \beta^2(\Gamma(1 + 1/\alpha))^2 $$
$$ E(X^2) = \beta^2\Gamma(1 + 2/\alpha) $$

Ahora igualamos los momentos poblacionales a los muestrales:
1. $E(X) = \bar{x} \implies \beta\Gamma(1 + 1/\alpha) = \bar{x}$  (Ecuación 1)
2. $E(X^2) = \frac{1}{n}\sum x_i^2 \implies \beta^2\Gamma(1 + 2/\alpha) = \frac{1}{n}\sum x_i^2$ (Ecuación 2)

Para encontrar una ecuación que implique solo la función gamma para estimar $\alpha$, dividimos el cuadrado de la Ecuación 1 entre la Ecuación 2:

$$ \frac{[E(X)]^2}{E(X^2)} = \frac{[\beta\Gamma(1 + 1/\alpha)]^2}{\beta^2\Gamma(1 + 2/\alpha)} = \frac{[\Gamma(1 + 1/\alpha)]^2}{\Gamma(1 + 2/\alpha)} $$

Sustituyendo los estimadores muestrales en el lado izquierdo:
$$ \frac{(\bar{x})^2}{\frac{1}{n}\sum x_i^2} = \frac{[\Gamma(1 + 1/\alpha)]^2}{\Gamma(1 + 2/\alpha)} $$

Esta es la ecuación que permite estimar $\alpha$. Una vez obtenido $\hat{\alpha}$, podemos despejar $\beta$ de la Ecuación 1:
$$ \hat{\beta} = \frac{\bar{x}}{\Gamma(1 + 1/\hat{\alpha})} $$

\textbf{Cálculo numérico:}

Datos: $n = 20$, $\bar{x} = 28.0$, $\sum x_i^2 = 16500$.

Calculamos el segundo momento muestral:
$$ \frac{1}{n}\sum x_i^2 = \frac{16500}{20} = 825 $$

Calculamos la razón del lado izquierdo de nuestra ecuación derivada:
$$ \text{Razón} = \frac{(28.0)^2}{825} = \frac{784}{825} \approx 0.9503 $$

Usando la ayuda proporcionada en el problema:
$$ \frac{[\Gamma(1.2)]^2}{\Gamma(1.4)} = 0.95 $$

Igualamos la estructura de la ecuación con el valor dado:
$$ \frac{[\Gamma(1 + 1/\alpha)]^2}{\Gamma(1 + 2/\alpha)} \approx 0.95 $$
Esto implica que:
$$ 1 + 1/\alpha = 1.2 \implies 1/\alpha = 0.2 \implies \alpha = 5 $$

Ahora estimamos $\beta$ usando $\hat{\alpha} = 5$ y la Ecuación 1:
$$ \bar{x} = \hat{\beta} \Gamma(1 + 1/5) = \hat{\beta} \Gamma(1.2) $$
$$ 28.0 = \hat{\beta} \Gamma(1.2) $$

Consultando tablas de la función Gamma o usando una calculadora, $\Gamma(1.2) \approx 0.9182$.
$$ \hat{\beta} = \frac{28.0}{0.9182} \approx 30.49 $$

\textbf{Resultados finales:}
$$ \hat{\alpha} = 5 $$
$$ \hat{\beta} \approx 30.49 $$


\section{Problema 5}

\subsection{Enunciado}
\textbf{Edades de estrellas.} Se calculan las edades de algunas estrellas dentro de cierto cúmulo de estrellas joven (en millones de años):
\[ 392, 376, 401, 367, 389, 362, 409, 415, 358, 375. \]

\textit{a)} Suponiendo que la edad de las estrellas en esta región sigue una distribución normal, calcule el promedio real y la desviación estándar usando el método de maximum likelihood.

\textit{b)} Una vez más, suponiendo una distribución normal, calcule la edad por debajo de la cual se encuentra el 95\% de todas las estrellas de este cúmulo. [Sugerencia: ¿cuál es el percentil 95 en términos de $\mu$ o $\sigma$? Ahora use el principio de invarianza.]

\textit{c)} Suponga que se decide medir la edad de las estrellas en otra región de este cúmulo. Sea $X$ dicha medida, use los datos dados para estimar la probabilidad de obtener un $X$ a lo sumo de 400 millones de años. [Sugerencia: $P(X \le 400) = \Phi((400 - \mu)/\sigma)$.]

\subsection{Solución}

Para resolver este problema, utilizaremos los conceptos de estimación de máxima verosimilitud (Maximum Likelihood Estimation, MLE) para una distribución normal, tal como se discute en el Capítulo 6 del libro Devore.

Sean los datos observados $x_1, x_2, \dots, x_{10}$:
\[ 392, 376, 401, 367, 389, 362, 409, 415, 358, 375 \]
El tamaño de la muestra es $n = 10$.

\subsubsection*{Inciso a)}
Se asume que la variable aleatoria $X$ (edad de las estrellas) sigue una distribución normal $N(\mu, \sigma^2)$. Debemos encontrar los estimadores de máxima verosimilitud para la media $\mu$ y la desviación estándar $\sigma$.

Según la Sección 6.2 del libro Devore, para una distribución normal, los estimadores de máxima verosimilitud son:
\begin{itemize}
    \item Para la media $\mu$: $\hat{\mu} = \bar{x} = \frac{1}{n}\sum_{i=1}^{n} x_i$
    \item Para la varianza $\sigma^2$: $\hat{\sigma}^2 = \frac{1}{n}\sum_{i=1}^{n} (x_i - \bar{x})^2$
\end{itemize}

\textbf{Cálculo de $\hat{\mu}$:}
\begin{align*}
    \sum_{i=1}^{10} x_i &= 392 + 376 + 401 + 367 + 389 + 362 + 409 + 415 + 358 + 375 = 3844 \\
    \hat{\mu} = \bar{x} &= \frac{3844}{10} = 384.4
\end{align*}
El estimado del promedio real es 384.4 millones de años.

\textbf{Cálculo de $\hat{\sigma}$:}
Primero calculamos la suma de los cuadrados de las desviaciones respecto a la media:
\begin{align*}
    \sum_{i=1}^{10} (x_i - \bar{x})^2 &= (392-384.4)^2 + (376-384.4)^2 + \dots + (375-384.4)^2 \\
    &= (7.6)^2 + (-8.4)^2 + (16.6)^2 + (-17.4)^2 + (4.6)^2 + (-22.4)^2 + (24.6)^2 + (30.6)^2 + (-26.4)^2 + (-9.4)^2 \\
    &= 57.76 + 70.56 + 275.56 + 302.76 + 21.16 + 501.76 + 605.16 + 936.36 + 696.96 + 88.36 \\
    &= 3556.4
\end{align*}
Ahora, el estimado de la varianza es:
\[ \hat{\sigma}^2 = \frac{3556.4}{10} = 355.64 \]
Y el estimado de la desviación estándar es:
\[ \hat{\sigma} = \sqrt{355.64} \approx 18.858 \]

\subsubsection*{Inciso b)}
Se pide calcular la edad por debajo de la cual se encuentra el 95\% de las estrellas. Esto corresponde al percentil 95 de la distribución, denotado como $\pi_{0.95}$.

Para una distribución normal $N(\mu, \sigma^2)$, el percentil $p$-ésimo está dado por:
\[ \pi_p = \mu + z_p \sigma \]
donde $z_p$ es el percentil $p$-ésimo de la distribución normal estándar $N(0,1)$.

Para $p = 0.95$, buscamos $z_{0.95}$ en la Tabla A.3 del apéndice. El valor que deja un área de 0.95 a su izquierda es aproximadamente $z_{0.95} = 1.645$.
Por lo tanto, el percentil 95 en términos de los parámetros es:
\[ \pi_{0.95} = \mu + 1.645\sigma \]

El \textbf{principio de invarianza} (Sección 6.2, Proposición) establece que si $\hat{\theta}_1, \dots, \hat{\theta}_k$ son los estimadores de máxima verosimilitud de los parámetros $\theta_1, \dots, \theta_k$, entonces el estimado de máxima verosimilitud de cualquier función $h(\theta_1, \dots, \theta_k)$ es $h(\hat{\theta}_1, \dots, \hat{\theta}_k)$.

Aplicando este principio, el estimado del percentil 95 es:
\[ \hat{\pi}_{0.95} = \hat{\mu} + 1.645\hat{\sigma} \]
Sustituyendo los valores calculados en el inciso a):
\begin{align*}
    \hat{\pi}_{0.95} &= 384.4 + 1.645(18.858) \\
    &\approx 384.4 + 31.021 \\
    &\approx 415.42
\end{align*}
La edad estimada por debajo de la cual se encuentra el 95\% de las estrellas es aproximadamente 415.42 millones de años.

\subsubsection*{Inciso c)}
Se desea estimar la probabilidad $P(X \le 400)$.
Estandarizando la variable $X \sim N(\mu, \sigma^2)$, tenemos:
\[ P(X \le 400) = P\left(\frac{X - \mu}{\sigma} \le \frac{400 - \mu}{\sigma}\right) = \Phi\left(\frac{400 - \mu}{\sigma}\right) \]
donde $\Phi$ es la función de distribución acumulativa de la normal estándar.

Nuevamente, por el principio de invarianza, estimamos esta probabilidad sustituyendo $\mu$ y $\sigma$ por sus estimadores de máxima verosimilitud $\hat{\mu}$ y $\hat{\sigma}$:
\begin{align*}
    \widehat{P(X \le 400)} &= \Phi\left(\frac{400 - \hat{\mu}}{\hat{\sigma}}\right) \\
    &= \Phi\left(\frac{400 - 384.4}{18.858}\right) \\
    &= \Phi\left(\frac{15.6}{18.858}\right) \\
    &\approx \Phi(0.827)
\end{align*}
Interpolando en la Tabla A.3 (Áreas de la Curva normal estándar):
\begin{itemize}
    \item $\Phi(0.82) = 0.7939$
    \item $\Phi(0.83) = 0.7967$
\end{itemize}
Para $z = 0.827$, el valor es aproximadamente $0.7939 + 0.7(0.7967 - 0.7939) \approx 0.7959$.

Por lo tanto, la probabilidad estimada de obtener una edad de a lo sumo 400 millones de años es aproximadamente \textbf{0.796}.


\section{Problema 6}

\subsection{Enunciado}

a) Sea $X$ una variable aleatoria que sigue una distribución binomial negativa. Muestre que para $r \geq 2$, un estimador no sesgado de la proporción $p$ está dado por:
$$\hat{p} = \frac{r-1}{X+r-1}$$

b) Encuentre el mle para $p$

c) Un vehículo con un defecto particular en su sistema de control de emisiones es llevado a diferentes mecánicos seleccionados al azar hasta que $r$ de ellos hayan diagnosticado correctamente el problema. Supongamos que esto requiere diagnósticos por parte de $n$ mecánicos diferentes (por lo que hubo $n-r$ diagnósticos incorrectos). Sea $p = P(\text{diagnóstico correcto})$, por lo que $p$ es la proporción de todos los mecánicos que diagnosticarían correctamente el problema. ¿Cuál es el mle de $p$? ¿Es lo mismo que el mle si una muestra aleatoria de 20 mecánicos da como resultado 3 diagnósticos correctos? Explique.

d) Compare el mle de $p$ con el estimador en (a).

\subsection{Solución}

\textbf{Inciso a)}

Para demostrar que $\hat{p} = \frac{r-1}{X+r-1}$ es un estimador insesgado de $p$, necesitamos mostrar que $E[\hat{p}] = p$.

La función de masa de probabilidad de la distribución binomial negativa es:
$$P(X = x) = \binom{x+r-1}{r-1} p^r (1-p)^x, \quad x = 0, 1, 2, \ldots$$

Calculamos el valor esperado:
$$E\left[\frac{r-1}{X+r-1}\right] = \sum_{x=0}^{\infty} \frac{r-1}{x+r-1} \binom{x+r-1}{r-1} p^r (1-p)^x$$

Notemos que:
$$\frac{r-1}{x+r-1} \binom{x+r-1}{r-1} = \frac{r-1}{x+r-1} \cdot \frac{(x+r-1)!}{(r-1)!x!} = \frac{(x+r-2)!}{(r-2)!x!} = \binom{x+r-2}{r-2}$$

Por lo tanto:
$$E\left[\frac{r-1}{X+r-1}\right] = \sum_{x=0}^{\infty} \binom{x+r-2}{r-2} p^r (1-p)^x$$

$$= p \sum_{x=0}^{\infty} \binom{x+r-2}{r-2} p^{r-1} (1-p)^x$$

La suma es igual a 1 porque es la suma de todas las probabilidades de una distribución binomial negativa con parámetros $r-1$ y $p$. Por lo tanto:

$$E\left[\frac{r-1}{X+r-1}\right] = p \cdot 1 = p$$

Esto demuestra que $\hat{p} = \frac{r-1}{X+r-1}$ es un estimador insesgado de $p$ para $r \geq 2$.

\textbf{Inciso b)}

Para encontrar el estimador de máxima verosimilitud (mle) de $p$, maximizamos la función de verosimilitud:

$$L(p) = \binom{x+r-1}{r-1} p^r (1-p)^x$$

La log-verosimilitud es:
$$\ell(p) = \ln L(p) = \ln\binom{x+r-1}{r-1} + r\ln p + x\ln(1-p)$$

Derivando con respecto a $p$ e igualando a cero:
$$\frac{d\ell}{dp} = \frac{r}{p} - \frac{x}{1-p} = 0$$

Resolviendo:
$$\frac{r}{p} = \frac{x}{1-p}$$
$$r(1-p) = xp$$
$$r - rp = xp$$
$$r = p(x+r)$$
$$\hat{p}_{\text{mle}} = \frac{r}{x+r}$$

\textbf{Inciso c)}

En el contexto del problema, $n$ es el número total de mecánicos consultados hasta obtener $r$ diagnósticos correctos. Esto significa que $X = n - r$ (número de diagnósticos incorrectos).

El mle de $p$ es:
$$\hat{p}_{\text{mle}} = \frac{r}{n-r+r} = \frac{r}{n}$$

Si una muestra aleatoria de 20 mecánicos da como resultado 3 diagnósticos correctos, el mle sería:
$$\hat{p}_{\text{mle}} = \frac{3}{20} = 0.15$$

Sin embargo, en el esquema binomial negativo, $n$ no es fijo de antemano; es el número de ensayos necesarios para obtener $r$ éxitos. En el esquema binomial (muestra fija de 20), el mle es $\frac{3}{20}$, pero en el esquema binomial negativo, si necesitamos $r=3$ éxitos y esto requiere $n$ ensayos, el mle es $\frac{3}{n}$.

No son lo mismo porque los esquemas de muestreo son diferentes: en uno el número de ensayos es fijo (binomial), en el otro el número de éxitos es fijo (binomial negativo).

\textbf{Inciso d)}

Comparando los estimadores:
- Estimador insesgado: $\hat{p} = \frac{r-1}{X+r-1} = \frac{r-1}{n-1}$
- MLE: $\hat{p}_{\text{mle}} = \frac{r}{X+r} = \frac{r}{n}$

El MLE tiende a sobreestimar ligeramente $p$ comparado con el estimador insesgado, ya que $\frac{r}{n} > \frac{r-1}{n-1}$ para $r < n$. El estimador insesgado corrige este sesgo restando 1 tanto en el numerador como en el denominador.

Para muestras grandes ($n$ grande), ambos estimadores son muy similares, pero para muestras pequeñas, el estimador insesgado proporciona una mejor estimación en el sentido de tener valor esperado igual al parámetro verdadero.

\section{Problema 7}

\subsection{Enunciado}
\textbf{Atmósferas planetarias} Un estudio de la composición de atmósferas en exoplanetas similares a la tierra encontró que para una muestra de 50 planetas, la media de la muestra del volumen de CO$_2$ fue 327.08 ppmv (partes por millón por volumen), y la desviación estándar de la muestra fue 82.215 ppmv.

\begin{itemize}
    \item[a)] Calcule e interprete un intervalo de confianza del 95\% para la cantidad promedio real de CO$_2$ en todos los planetas como la Tierra.
    \item[b)] Suponga que los investigadores tomaron un valor aproximado de 87.5 para $s$ antes de recopilar los datos. ¿Qué tamaño de muestra sería necesario para obtener un ancho de intervalo de 25 ppmv para un nivel de confianza del 95\%?
\end{itemize}

\subsection{Solución}

\subsubsection{Inciso a)}
Para resolver este problema, primero identificamos los datos proporcionados en el enunciado:
\begin{itemize}
    \item Tamaño de la muestra: $n = 50$
    \item Media muestral: $\bar{x} = 327.08$ ppmv
    \item Desviación estándar muestral: $s = 82.215$ ppmv
    \item Nivel de confianza: $95\%$
\end{itemize}

Dado que el tamaño de la muestra es grande ($n \geq 30$), es apropiado utilizar la distribución normal estándar ($z$) para construir el intervalo de confianza para la media poblacional $\mu$, incluso si la desviación estándar poblacional $\sigma$ es desconocida y estimada por $s$. La fórmula general para un intervalo de confianza de gran muestra es:

$$ \bar{x} \pm z_{\alpha/2} \left( \frac{s}{\sqrt{n}} \right) $$

Donde:
\begin{itemize}
    \item $\bar{x}$ es la media muestral.
    \item $z_{\alpha/2}$ es el valor crítico de la distribución normal estándar que deja un área de $\alpha/2$ en la cola superior.
    \item $s$ es la desviación estándar muestral.
    \item $n$ es el tamaño de la muestra.
\end{itemize}

Para un nivel de confianza del 95\%, tenemos que $1 - \alpha = 0.95$, por lo que $\alpha = 0.05$. Esto significa que $\alpha/2 = 0.025$. Buscando en la tabla de la distribución normal estándar (Tabla A.3 del apéndice del libro Devore), el valor $z$ que corresponde a un área acumulada de $1 - 0.025 = 0.975$ es $z_{0.025} = 1.96$.

Sustituimos los valores en la fórmula:

$$ 327.08 \pm 1.96 \left( \frac{82.215}{\sqrt{50}} \right) $$

Primero calculamos el error estándar de la media ($s/\sqrt{n}$):

$$ \frac{82.215}{\sqrt{50}} \approx \frac{82.215}{7.07106} \approx 11.627 $$

Luego calculamos el margen de error (la cantidad que se suma y resta):

$$ 1.96 \times 11.627 \approx 22.789 $$

Finalmente, construimos el intervalo sumando y restando el margen de error a la media muestral:

$$ \text{Límite inferior} = 327.08 - 22.79 = 304.29 $$
$$ \text{Límite superior} = 327.08 + 22.79 = 349.87 $$

El intervalo de confianza del 95\% es $(304.29, 349.87)$.

\textbf{Interpretación:} Podemos afirmar con un 95\% de confianza que el intervalo $(304.29, 349.87)$ contiene la verdadera media poblacional del volumen de CO$_2$ en todos los exoplanetas similares a la Tierra. Esto significa que si tomáramos muchas muestras de 50 planetas y construyéramos un intervalo de confianza para cada una, aproximadamente el 95\% de esos intervalos contendrían la media real $\mu$.

\subsubsection{Inciso b)}
En esta parte, se nos pide determinar el tamaño de muestra $n$ necesario para lograr una precisión específica en la estimación. Los datos son:
\begin{itemize}
    \item Ancho del intervalo deseado ($w$) = 25 ppmv.
    \item Valor estimado de la desviación estándar ($s$) = 87.5 ppmv (valor de planificación).
    \item Nivel de confianza = 95\%, por lo que $z_{\alpha/2} = 1.96$.
\end{itemize}

El ancho de un intervalo de confianza está dado por $2 \times (\text{margen de error})$. Es decir:
$$ w = 2 z_{\alpha/2} \frac{s}{\sqrt{n}} $$

Queremos despejar $n$ de esta ecuación. Primero dividimos por 2 a ambos lados para obtener el margen de error ($E = w/2 = 12.5$):
$$ \frac{w}{2} = z_{\alpha/2} \frac{s}{\sqrt{n}} $$
$$ 12.5 = 1.96 \frac{87.5}{\sqrt{n}} $$

Despejamos $\sqrt{n}$:
$$ \sqrt{n} = \frac{1.96 \times 87.5}{12.5} $$

Calculamos el numerador:
$$ 1.96 \times 87.5 = 171.5 $$

Dividimos por el denominador:
$$ \sqrt{n} = \frac{171.5}{12.5} = 13.72 $$

Elevamos al cuadrado para encontrar $n$:
$$ n = (13.72)^2 $$
$$ n \approx 188.2384 $$

Dado que el tamaño de la muestra debe ser un número entero de planetas, y para garantizar que el ancho del intervalo sea \textit{al menos} tan pequeño como el deseado (o que la precisión sea al menos la requerida), siempre redondeamos hacia arriba al siguiente número entero.

$$ n = 189 $$

Por lo tanto, sería necesario una muestra de \textbf{189 planetas} para obtener un intervalo de confianza del 95\% con un ancho de 25 ppmv, asumiendo una desviación estándar de 87.5 ppmv.



\section{Problema 8}

\subsection{Enunciado}
\textbf{Binarias eclipsantes.} Algunas veces, bajas en la curva de luz de una estrella es atribuida a un tránsito planetario. Sin embargo, este no es el único fenómeno que produce estos cambios en dichas curvas. En ocasiones estos bajones son debido a una estrella compañera pasando ya sea por detrás o por delante de la otra. A este fenómeno se le conoce como binaria eclipsante. Suponga que en una muestra de 88 curvas reportadas como tránsito, se encontró que 31 eran binarias eclipsantes.

\begin{enumerate}
    \item[a)] Calcule e interprete un intervalo de confianza (IC) usando un nivel de confianza del 95\% para la proporción real de todas las curvas de luz encontradas con binarias eclipsantes bajo estas circunstancias.
    \item[b)] ¿Qué tamaño de muestra se requiere si el ancho deseado del IC del 95\% debe ser como máximo .04, independientemente de los resultados de la muestra?
    \item[c)] Se desea conocer solamente el límite de confianza superior del 95\%, en lugar del intervalo. ¿El límite superior del intervalo en (a) corresponde al límite deseado para la proporción estimada? Explique.
\end{enumerate}

\subsection{Solución}

Este problema se basa en la teoría de intervalos de confianza para una proporción de población, tratada en la \textbf{Sección 7.2} del libro de texto.

\subsubsection{Inciso a)}
Primero, identificamos los datos proporcionados:
\begin{itemize}
    \item Tamaño de la muestra: $n = 88$
    \item Número de éxitos (binarias eclipsantes): $x = 31$
    \item Nivel de confianza: $1 - \alpha = 0.95$, por lo tanto $\alpha = 0.05$.
\end{itemize}

Calculamos la proporción muestral $\hat{p}$:
$$ \hat{p} = \frac{x}{n} = \frac{31}{88} \approx 0.3523 $$

Para un nivel de confianza del 95\%, el valor crítico $z_{\alpha/2}$ es $z_{0.025}$. De la tabla de la distribución normal estándar (Tabla A.3), sabemos que $z_{0.025} = 1.96$.

El intervalo de confianza de muestra grande para una proporción $p$ está dado por la fórmula:
$$ \hat{p} \pm z_{\alpha/2} \sqrt{\frac{\hat{p}(1-\hat{p})}{n}} $$

Sustituyendo los valores:
$$ 0.3523 \pm 1.96 \sqrt{\frac{0.3523(1-0.3523)}{88}} $$
$$ 0.3523 \pm 1.96 \sqrt{\frac{0.3523(0.6477)}{88}} $$
$$ 0.3523 \pm 1.96 \sqrt{\frac{0.2282}{88}} $$
$$ 0.3523 \pm 1.96 \sqrt{0.002593} $$
$$ 0.3523 \pm 1.96 (0.0509) $$
$$ 0.3523 \pm 0.0998 $$

El intervalo de confianza es $(0.3523 - 0.0998, 0.3523 + 0.0998)$, es decir:
$$ (0.2525, 0.4521) $$

\textbf{Interpretación:} Con un 95\% de confianza, podemos afirmar que la proporción real de todas las curvas de luz reportadas como tránsito que en realidad son binarias eclipsantes se encuentra entre 0.2525 y 0.4521.

\subsubsection{Inciso b)}
Se desea un ancho de intervalo $w = 0.04$. El ancho del intervalo es el doble del margen de error ($w = 2E$), por lo que el margen de error deseado es $E = 0.02$.

La fórmula para el tamaño de muestra $n$ necesario para estimar $p$ con un margen de error $E$ es:
$$ n = \frac{z_{\alpha/2}^2 \hat{p}(1-\hat{p})}{E^2} $$

Sin embargo, el problema pide que el tamaño de muestra sea válido "independientemente de los resultados de la muestra". Esto significa que debemos usar el valor conservador para $\hat{p}$ que maximiza el producto $\hat{p}(1-\hat{p})$. Este valor ocurre cuando $\hat{p} = 0.5$.

Sustituyendo $z_{0.025} = 1.96$, $\hat{p} = 0.5$ y $E = 0.02$:
$$ n = \frac{(1.96)^2 (0.5)(0.5)}{(0.02)^2} $$
$$ n = \frac{3.8416 \cdot 0.25}{0.0004} $$
$$ n = \frac{0.9604}{0.0004} $$
$$ n = 2401 $$

Se requiere un tamaño de muestra de al menos \textbf{2401} curvas de luz.

\subsubsection{Inciso c)}
No, el límite superior del intervalo en (a) no corresponde al límite de confianza superior unilateral del 95\% deseado.

\textbf{Explicación:}
\begin{itemize}
    \item El intervalo calculado en el inciso (a) es un intervalo de confianza \textbf{bilateral} (de dos colas) del 95\%. Esto significa que el 5\% de error ($\alpha$) se divide en dos colas: 2.5\% en la cola inferior y 2.5\% en la cola superior. Por lo tanto, el límite superior se calcula usando el valor crítico $z_{\alpha/2} = z_{0.025} = 1.96$. Este límite superior deja un 2.5\% de probabilidad en la cola superior, lo que corresponde a un nivel de confianza del 97.5\% para un límite superior unilateral.
    \item Un límite de confianza superior \textbf{unilateral} del 95\% coloca todo el 5\% de error ($\alpha$) en la cola superior. Por lo tanto, se calcularía usando el valor crítico $z_{\alpha} = z_{0.05} = 1.645$.
\end{itemize}

Dado que $1.96 \neq 1.645$, los límites son diferentes. El límite superior del intervalo bilateral en (a) será más alto (más conservador) que el límite superior unilateral del 95\%.


\section{Problema 9}
\subsection{Enunciado}
Esperando el bus. Sea $X_1, X_2, ..., X_n$ una muestra aleatoria producto de una distribución uniforme en el intervalo $[0, \theta]$, de la forma:
$$
f(x) = \begin{cases} \frac{1}{\theta} & 0 \le x \le \theta \\ 0 & \text{en otro caso} \end{cases}
$$
entonces si $Y = \max(X_i)$, se puede mostrar que la variable aleatoria $U = Y/\theta$ tiene una función de densidad dada por:
$$
f_U(u) = \begin{cases} nu^{n-1} & 0 \le u \le 1 \\ 0 & \text{en otro caso} \end{cases}
$$
\begin{enumerate}
    \item[a)] Use $f_U(u)$ para verificar que:
    $$ P\left( (\alpha/2)^{1/n} < \frac{Y}{\theta} \le (1 - \alpha/2)^{1/n} \right) = 1 - \alpha $$
    y use esto para derivar el IC $100(1-\alpha)\%$ para $\theta$.
    \item[b)] Verifique que $P(\alpha^{1/n} \le Y/\theta \le 1) = 1 - \alpha$, y derive un IC $100(1-\alpha)\%$ para $\theta$ basado en dicha probabilidad.
    \item[c)] ¿Cuál de los dos intervalos derivados anteriormente es más corto? Si el tiempo de espera de un bus para ir al centro se distribuye uniformemente y los tiempos de espera observados son $x_1 = 4.2, x_2 = 3.5, x_3 = 1.7, x_4 = 1.2$ y $x_5 = 2.4$, obtenga un IC del 95\% para $\theta$ utilizando el intervalo más corto.
\end{enumerate}

\subsection{Solución}

\subsubsection{Inciso a)}
Para verificar la probabilidad, integramos la función de densidad de $U$ en el intervalo dado:
$$
\begin{aligned}
P\left( (\alpha/2)^{1/n} < U \le (1 - \alpha/2)^{1/n} \right) &= \int_{(\alpha/2)^{1/n}}^{(1 - \alpha/2)^{1/n}} f_U(u) \, du \\
&= \int_{(\alpha/2)^{1/n}}^{(1 - \alpha/2)^{1/n}} n u^{n-1} \, du \\
&= \left[ u^n \right]_{(\alpha/2)^{1/n}}^{(1 - \alpha/2)^{1/n}} \\
&= \left( (1 - \alpha/2)^{1/n} \right)^n - \left( (\alpha/2)^{1/n} \right)^n \\
&= (1 - \alpha/2) - (\alpha/2) \\
&= 1 - \alpha
\end{aligned}
$$
La verificación es correcta. Ahora, para derivar el intervalo de confianza, partimos de la desigualdad de probabilidad:
$$ (\alpha/2)^{1/n} < \frac{Y}{\theta} \le (1 - \alpha/2)^{1/n} $$
Invertimos los términos (notando que todos son positivos):
$$ \frac{1}{(1 - \alpha/2)^{1/n}} \le \frac{\theta}{Y} < \frac{1}{(\alpha/2)^{1/n}} $$
Multiplicamos por $Y$ para aislar $\theta$:
$$ \frac{Y}{(1 - \alpha/2)^{1/n}} \le \theta < \frac{Y}{(\alpha/2)^{1/n}} $$
Por lo tanto, un intervalo de confianza del $100(1-\alpha)\%$ para $\theta$ es:
$$ \left[ \frac{Y}{(1 - \alpha/2)^{1/n}}, \frac{Y}{(\alpha/2)^{1/n}} \right) $$

\subsubsection{Inciso b)}
Verificamos la segunda probabilidad de manera similar:
$$
\begin{aligned}
P\left( \alpha^{1/n} \le U \le 1 \right) &= \int_{\alpha^{1/n}}^{1} n u^{n-1} \, du \\
&= \left[ u^n \right]_{\alpha^{1/n}}^{1} \\
&= 1^n - (\alpha^{1/n})^n \\
&= 1 - \alpha
\end{aligned}
$$
La verificación es correcta. Derivamos el intervalo de confianza partiendo de:
$$ \alpha^{1/n} \le \frac{Y}{\theta} \le 1 $$
Invertimos los términos:
$$ 1 \le \frac{\theta}{Y} \le \frac{1}{\alpha^{1/n}} $$
Multiplicamos por $Y$:
$$ Y \le \theta \le \frac{Y}{\alpha^{1/n}} $$
Por lo tanto, un segundo intervalo de confianza del $100(1-\alpha)\%$ para $\theta$ es:
$$ \left[ Y, \frac{Y}{\alpha^{1/n}} \right] $$

\subsubsection{Inciso c)}
Para determinar cuál intervalo es más corto, comparamos sus longitudes. Sea $L_a$ la longitud del intervalo del inciso (a) y $L_b$ la del inciso (b).
$$ L_a = \frac{Y}{(\alpha/2)^{1/n}} - \frac{Y}{(1 - \alpha/2)^{1/n}} = Y \left[ \left(\frac{2}{\alpha}\right)^{1/n} - \left(\frac{1}{1 - \alpha/2}\right)^{1/n} \right] $$
$$ L_b = \frac{Y}{\alpha^{1/n}} - Y = Y \left[ \left(\frac{1}{\alpha}\right)^{1/n} - 1 \right] $$
Dado que $\alpha$ es pequeño (típicamente 0.05 o 0.01), el término $(2/\alpha)^{1/n}$ en $L_a$ es significativamente mayor que $(1/\alpha)^{1/n}$ en $L_b$. Además, restamos un valor cercano a 1 en $L_a$, mientras que en $L_b$ restamos exactamente 1. Numéricamente, para valores típicos de $\alpha$ y $n$, se puede demostrar que $L_b < L_a$. Por lo tanto, el intervalo del inciso (b) es más corto.

Ahora, calculamos el intervalo del 95\% utilizando el método del inciso (b).
Los datos son: $x_1 = 4.2, x_2 = 3.5, x_3 = 1.7, x_4 = 1.2, x_5 = 2.4$.
El tamaño de la muestra es $n = 5$.
El estadístico $Y$ es el máximo de las observaciones: $Y = \max(4.2, 3.5, 1.7, 1.2, 2.4) = 4.2$.
Para un IC del 95\%, tenemos $1 - \alpha = 0.95$, por lo que $\alpha = 0.05$.

Sustituimos en la fórmula del intervalo (b):
$$ \left[ 4.2, \frac{4.2}{(0.05)^{1/5}} \right] $$
Calculamos el límite superior:
$$ (0.05)^{1/5} = (0.05)^{0.2} \approx 0.54928 $$
$$ \text{Límite Superior} = \frac{4.2}{0.54928} \approx 7.646 $$

El intervalo de confianza del 95\% para $\theta$ es aproximadamente \textbf{[4.2, 7.65]}.

\textit{Referencia: Sección 7.1 del libro (Propiedades básicas de los intervalos de confianza) y conceptos de estimación por intervalos basados en estadísticos suficientes/pivote (temas que se extienden en capítulos posteriores sobre teoría de estimación).}


\section{Problema 10}

\textbf{Enunciado:}

\textit{Velocidad de dispersión.} Sea $\sigma$ la velocidad de dispersión estelar de una galaxia, medida en km/s. Suponga que el valor de $\log(\sigma)$ de las muestras de galaxias tomadas de cualquier cúmulo de galaxias particular se distribuye normalmente con una desviación estándar real de 0.4.

\begin{enumerate}
    \item[a)] Calcule un intervalo de confianza (IC) del 95\% para el valor de $\log(\sigma)$ promedio real de un cúmulo de galaxias determinado si su valor promedio para una muestra de 20 galaxias fue de 2.42.
    \item[b)] Calcule un IC del 98\% para $\log(\sigma)$ promedio real de otro cúmulo, basado en 16 elementos de una muestra con un valor de $\log(\sigma)$ promedio de 2.28.
    \item[c)] ¿Qué tamaño de muestra es necesario para que el ancho del intervalo del 95\% sea 0.2?
    \item[d)] ¿Qué tamaño de muestra es necesario para estimar el promedio real de $\log(\sigma)$ dentro de 0.1 del valor muestral, con un 99\% de confianza?
\end{enumerate}

\subsection{Solución}

Este problema trata sobre la construcción de intervalos de confianza para la media poblacional ($\mu$) cuando la desviación estándar poblacional ($\sigma$) es conocida y la población se distribuye normalmente. Utilizaremos la distribución normal estándar ($Z$).

La fórmula general para un intervalo de confianza es:
$$ \bar{x} \pm z_{\alpha/2} \frac{\sigma}{\sqrt{n}} $$
Donde:
\begin{itemize}
    \item $\bar{x}$ es la media muestral.
    \item $\sigma$ es la desviación estándar poblacional (dada como 0.4).
    \item $n$ es el tamaño de la muestra.
    \item $z_{\alpha/2}$ es el valor crítico de la distribución normal estándar que deja un área de $\alpha/2$ a su derecha.
\end{itemize}

\subsubsection{Inciso a)}
\textbf{Datos:}
\begin{itemize}
    \item Nivel de confianza: $95\% \Rightarrow 1 - \alpha = 0.95 \Rightarrow \alpha = 0.05$.
    \item Valor crítico: $z_{\alpha/2} = z_{0.025} = 1.96$.
    \item Tamaño de muestra: $n = 20$.
    \item Media muestral: $\bar{x} = 2.42$.
    \item Desviación estándar: $\sigma = 0.4$.
\end{itemize}

Sustituyendo en la fórmula:
$$ IC_{95\%} = 2.42 \pm 1.96 \left( \frac{0.4}{\sqrt{20}} \right) $$
$$ IC_{95\%} = 2.42 \pm 1.96 \left( \frac{0.4}{4.4721} \right) $$
$$ IC_{95\%} = 2.42 \pm 1.96 (0.08944) $$
$$ IC_{95\%} = 2.42 \pm 0.1753 $$

El intervalo de confianza es $(2.2447, 2.5953)$.

\subsubsection{Inciso b)}
\textbf{Datos:}
\begin{itemize}
    \item Nivel de confianza: $98\% \Rightarrow 1 - \alpha = 0.98 \Rightarrow \alpha = 0.02$.
    \item Valor crítico: $z_{\alpha/2} = z_{0.01} \approx 2.326$ (o 2.33 según la tabla). Usaremos 2.326 para mayor precisión.
    \item Tamaño de muestra: $n = 16$.
    \item Media muestral: $\bar{x} = 2.28$.
    \item Desviación estándar: $\sigma = 0.4$.
\end{itemize}

Sustituyendo en la fórmula:
$$ IC_{98\%} = 2.28 \pm 2.326 \left( \frac{0.4}{\sqrt{16}} \right) $$
$$ IC_{98\%} = 2.28 \pm 2.326 \left( \frac{0.4}{4} \right) $$
$$ IC_{98\%} = 2.28 \pm 2.326 (0.1) $$
$$ IC_{98\%} = 2.28 \pm 0.2326 $$

El intervalo de confianza es $(2.0474, 2.5126)$.

\subsubsection{Inciso c)}
Se requiere que el ancho total del intervalo ($w$) sea 0.2. El ancho de un intervalo de confianza está dado por $w = 2 \cdot E$, donde $E$ es el margen de error ($z_{\alpha/2} \frac{\sigma}{\sqrt{n}}$).

\textbf{Datos:}
\begin{itemize}
    \item Ancho deseado ($w$) = 0.2.
    \item Nivel de confianza: $95\% \Rightarrow z_{0.025} = 1.96$.
    \item $\sigma = 0.4$.
\end{itemize}

La fórmula para el tamaño de muestra basada en el ancho es:
$$ n = \left( \frac{2 \cdot z_{\alpha/2} \cdot \sigma}{w} \right)^2 $$
$$ n = \left( \frac{2 \cdot 1.96 \cdot 0.4}{0.2} \right)^2 $$
$$ n = \left( \frac{1.568}{0.2} \right)^2 $$
$$ n = (7.84)^2 $$
$$ n = 61.4656 $$

Como el tamaño de muestra debe ser un número entero, redondeamos hacia arriba.
\textbf{Respuesta:} Se necesita una muestra de tamaño \textbf{62}.

\subsubsection{Inciso d)}
Se requiere estimar la media real dentro de 0.1 del valor muestral. Esto significa que el margen de error ($E$) debe ser 0.1.

\textbf{Datos:}
\begin{itemize}
    \item Margen de error ($E$) = 0.1.
    \item Nivel de confianza: $99\% \Rightarrow 1 - \alpha = 0.99 \Rightarrow \alpha = 0.01$.
    \item Valor crítico: $z_{\alpha/2} = z_{0.005} \approx 2.576$ (o 2.58). Usaremos 2.576.
    \item $\sigma = 0.4$.
\end{itemize}

La fórmula para el tamaño de muestra basada en el margen de error es:
$$ n = \left( \frac{z_{\alpha/2} \cdot \sigma}{E} \right)^2 $$
$$ n = \left( \frac{2.576 \cdot 0.4}{0.1} \right)^2 $$
$$ n = \left( \frac{1.0304}{0.1} \right)^2 $$
$$ n = (10.304)^2 $$
$$ n \approx 106.17 $$

Redondeando hacia arriba al siguiente entero.
\textbf{Respuesta:} Se necesita una muestra de tamaño \textbf{107}.

\vspace{1em}
\noindent \textbf{Referencia:} Devore, J. L. (2008). \textit{Probabilidad y Estadística para Ingeniería y Ciencias}. Séptima edición. Capítulo 7: Intervalos estadísticos basados en una sola muestra (Secciones 7.2 y 7.3).


\end{document}